\documentclass{article}
\usepackage{geometry}
\usepackage{concmath}
\usepackage[OT1]{fontenc}
\usepackage{graphicx}
\usepackage{amsfonts}
\usepackage{amsmath,mathtools}
\usepackage{calrsfs}
\usepackage{color}
\usepackage[dvipsnames]{xcolor}
\usepackage[many]{tcolorbox}
\usepackage{setspace}
\usepackage{tkz-tab}    
\usepackage{amssymb}          
\usepackage{multicol}  
\usepackage{esint}
\geometry{margin=0.5in}
\definecolor{green}{RGB}{52, 235, 168}
\definecolor{violet}{RGB}{165, 83, 237}
\definecolor{lemon}{RGB}{170, 237, 83}
\definecolor{red}{RGB}{240, 125, 84}
\definecolor{blue}{RGB}{105, 255, 242}

\begin{document}
\noindent \colorbox{green}{531.13}
Soit \(x \in [0,1[\) fixé\\
\begin{align*}
  f_{n}(x) &= \frac{1}{1 + x + \dots + x^{n}}\\
           &= \frac{1-x}{1-x^{n}}\\
           &\xrightarrow[n \to +\infty]{} 1-x
\end{align*}
Donc \(f_{n} \rightarrow (g:x \mapsto 1-x)\)\\
\underline{Dominons :}\\
\begin{align*}
  |f_{n}(x)| &= \frac{1}{1+x+\dots+x^{n}}\\
             &\leq 1 \text{ indépendant de n, intégrable sur }[0, 1[
\end{align*}
Par convergence dominée : 
\begin{align*}
  \int_{0}^{1} \frac{dx}{1+x+\dots+x^{n}} &\rightarrow \int_{0}^{1}1-x dx \\
    &= \frac{1}{2}
\end{align*}

\noindent \colorbox{green}{531.1}\\
\begin{eqnarray*}
    \int_{0}^{+\infty} \frac{\sin nt }{nt + t^{2}}dt\\
    f_{n}(t) = \frac{\sin nt}{nt + t^{2}}
\end{eqnarray*}
Soit \(t \in ]0, +\infty[\) fixé, 
\begin{align*}
  f_{n}(t) &= \frac{O(1)}{nt + t^{2}}\\
    &= O(\frac{1}{n})\\ 
    &\xrightarrow[n \to +\infty]{} 0
\end{align*}
\underline{Dominons}\\
Pour \(t \in ]0, +\infty[, n  \in \mathbb{N}\)
\begin{align*}
  |f_{n}(t)| &= \frac{|\sin nt|}{nt + t^{2}} \\
             &\leq \frac{1}{nt+t^{2}}\\
             &\leq \frac{1}{t^{2}} \text { indépendant de n, 
intégrable en \(+ \infty\) mais sur \(]0, +\infty[\)} &= \frac{1}{1+\frac{t}{n}} \\
             & &\leq 1 \\
             & &\text{indépendant de n, intégrable en 0, mais pas sur } ]0, +\infty[
\end{align*}
\begin{align*}
    \varphi(t) : t \mapsto \begin{cases} 
      1 \text{ si } t < 1 \\
      \frac{1}{t^{2}} \text{ si } t \geq 1
    \end{cases}\\
    |f_{n}(t)| \leq \varphi(t) \text{ indépendant de n}
\end{align*}
\(\varphi\) est continue par morceaux sur \(]0, +\infty[\)\\
Au voisinnage de 0 : \\
\(\varphi(t) = 1\) intégrable\\
Au voisinnage de \(+\infty\)\\
\(\varphi(t) = \frac{1}{t^{2}}\) intégrable\\
Donc \(\varphi\) est intégrable sur \(]0, +\infty[\)\\
Donc \(\int_{0}^{+\infty}f_{n}(t)dt \xrightarrow[n \to +\infty]{} \int_{0}^{+\infty}0 dt = 0\)\\

\noindent \colorbox{green}{531.9}\\
\begin{align*}
  f:[0,1] \rightarrow \mathbb{R}, f \in \mathcal{C}^{0}([0,1])\\ 
  \int_{0}^{1} f(t) dt \\
  g : t \mapsto f(t^{n})
\end{align*}
Convergence simple : \\
Pour \(t \in [0,1[\) fixé\\ 
\begin{align*}
  g_{n}(t) &= f(t^{n})\\
           &\rightarrow f(0) \text{ par continuité de f}
\end{align*}
Donc \(g_{n}\xrightarrow[n \to +\infty]{} (t \mapsto f(0))\)\\
Dominons \\
\begin{align*}
    |g_{n}(t)| = |f_{n}(t)| \\
    &\leq ||f||_{\infty}^{[0,1]} \text{ existe car f continue sur [0,1], indépendant de n, intégrable sur segment [0,1]}
\end{align*}
Par convergence dominée\\
\begin{equation*}
    \int_{0}^{1} g_{n}(t)dt \xrightarrow[n \rightarrow +\infty]{} \int_{0}^{1}f(0)dt = f(0)
\end{equation*}

\noindent \colorbox{green}{531.16}\\
On note \(f_{n}(x) = \frac{1}{(1+x^{3})^{n}}\)\\
Soit \(x \in ]0, +\infty[\) fixé
\begin{align*}
   f_{n}(x)  \xrightarrow[n \rightarrow +\infty]{} 0
\end{align*}
Donc \(f_{n} \xrightarrow[n \rightarrow +\infty]{cs} (x \mapsto 0)\)\\
Dominons \\
\begin{align*}
  |f_{n}(t)| \leq \frac{1}{1+x^{3}} \text{ indépendant de n, intégrable sur } ]0, +\infty[
\end{align*}
Donc par convergence dominée
\begin{align*}
  v_{n} &= \int_{0}^{+\infty}\frac{dt }{(1+t^{3})^{n}}\\
    v_{n}  \xrightarrow[n \rightarrow +\infty]{} &\int_{0}^{+\infty} 0 dt \\
    &= 0
    u_{n} &= (-1)^{n}v_{n}
\end{align*}
or \(v_{n}\) est décroissante car \((1 + t^{3})^{n} \leq (1+t^{3})^{n+1}\)

\noindent \colorbox{green}{531.20}\\
a/ 
\begin{align*}
    g_{n}(x) = \frac{f(x)}{1+x^{n}} \rightarrow \begin{cases}
      f(x) \text{ si } x \in [0, 1[ \\
      \frac{f(1)}{2} \text{ si x = 1}\\
      0 \text{ si } x \in ]1, b]
    \end{cases}
\end{align*}
Dominons : \(|g_{n}(x)| \leq |f(x)|\) indépendant de n \\
\(f(x)\) continue sur [a,b] segment donc intégrable\\
b/ Calculons par parties avec \(xf(x) \in \mathcal{C}^{1}\)
\begin{align*}
  \int_{0}^{1} \frac{x^{n-1}}{1+x^{n}}x f(x) dx &= [\frac{1}{n}ln(1+x^{n}xf(x))]^{1}_{a}-\int_{a}^{1}\frac{1}{n}ln(1+x^{n})(f(x)+xf'(x))dx\\
                                                &= \frac{1}{n}ln(2)f(1) - \frac{1}{n}ln(1+a^{n})af(a)-\frac{1}{n}J_{n}\\
                                                &= \frac{1}{n}ln(2)f(1) + \frac{1}{n}o(1) -\frac{1}{n}J_{n}
\end{align*}
Montrons que \(J_{n}\xrightarrow[n \to +\infty]{} 0\)
\begin{align*}
  f_{n}(x) &= ln(1+x^{n})(f(x)-xf'(x))
\end{align*}
Soit \(x \in [a, 1[\) fixé
\begin{align*}
  f_{n}(x) \xrightarrow[n\to +\infty]{} 0\\
  f_{n} \xrightarrow[]{cs} g:x \mapsto 0
\end{align*}
\underline{Dominons}
\begin{align*}
  |f_{n}(x)| &= ln(1+x^{n})|f(x) + xf'(x)| \\
             &\leq ln(2)|f(x)+xf'(x)|\\
             &\hookleftarrow \text{indé de n et continue sur [a,1] segment donc intégrable sur } [a,1[
\end{align*}
Par convergence dominée 
\begin{equation*}
  J_{n} \xrightarrow[n \to +\infty]{} \int_{a}^{1}g(x)dx = 0
\end{equation*}
\begin{align*}
  \text{Ainsi } I_{n}&=\frac{1}{n}ln(2)f(1) + \frac{1}{n}o(1)\\
                     &\backsim_{+\infty} \frac{1}{n}ln(2)f(1)
\end{align*}

















\end{document}

